\subsection{Validation-Selected False-Alarm Correction}
The balanced-loss reference was followed by an exploratory comparison of unweighted and square-root inverse-frequency weighting. Architecture, preprocessing, optimizer and the 200-epoch endpoint were held fixed. Candidate training used seed 42 and training/validation arrays only. Neither candidate met the initial 1\% validation false-alarm target with at least 90\% attack detection. Before generating new test predictions, the false-alarm limit was revised to 3\%, retaining the detection constraint. Square-root weighting was selected by validation macro-F1 (0.8552 versus 0.8480 for unweighted loss).

The exploratory validation-selected square-root weighting correction reduced mean normal test false alarms from 42.70\% to 1.90\%, with accuracy 95.44\%, macro-F1 0.8551 and attack detection 90.96\%. However, NMRI recall fell from 94.39\% to 0.54\%. This is a false-alarm versus attack-sensitivity trade-off, not a complete detector-quality solution; the original 1\% false-alarm target remains unmet. Independent evaluation is still required.
The selected setting was evaluated at paired final seeds 10042, 10043 and 10044, each trained for exactly 200 epochs on GPU. These results describe a conventional loss correction on the historical split. The revised target is exploratory, and no test-driven threshold or endpoint selection was performed after these new predictions.
